\section{Patched conics}

\subsection{Hipérbola geocéntrica}

Cuando Apophis entra en la esfera de influencia de la Tierra (aproximadamente $0.9$ millones de kilómetros), calcular los elementos orbitales de la hipérbola geocéntrica que describe su trayectoria respecto a la Tierra. Verificar que $e > 1$ y calcular el parámetro de impacto $b$ y el ángulo de deflexión.

\subsection{Comparación patched conics vs. N cuerpos}

Implementar la aproximación completa de \textit{patched conics} ---órbita heliocéntrica más hipérbola geocéntrica--- y comparar la distancia mínima predicha con la obtenida en la simulación de $N$ cuerpos. Cuantificar cuánto mejora esta predicción frente a la aproximación simple de 2 cuerpos.

\subsection{Visualización del cambio de cónica}

Graficar la trayectoria completa de Apophis con colores distintos según la región dinámica dominante: azul cuando domina el Sol (trayectoria heliocéntrica) y rojo cuando domina la Tierra (trayectoria geocéntrica). Marcar visualmente la esfera de influencia terrestre.